% Part: computability
% Chapter: computability-theory
% Section: complete-ce-sets

\documentclass[../../../include/open-logic-section]{subfiles}

\begin{document}

% SEGMENT OLP-0245-S01

\olfileid{cmp}{thy}{cce}
\olsection{முழுமையான கணிக்கத்தக்கவாறு எண்ணிடத்தக்க கணங்கள்}

\begin{defn}
பின்வருவன நிறைவேறினால், கணம் $A$ ஆனது (பலவற்றுக்கு ஒன்று
குறைத்தன்மையின்கீழ்) ஒரு
\emph{முழுமையான !!{computably enumerable} கணம்}:
\begin{enumerate}
\item $A$ கணிக்கத்தக்கவாறு எண்ணிடத்தக்கது; மேலும்
\item வேறு எந்தக் கணிக்கத்தக்கவாறு எண்ணிடத்தக்க கணம் $B$ க்கும்,
  $B \leq_m A$.
\end{enumerate}
\end{defn}

வேறுவிதமாகக் கூறினால், முழுமையான கணிக்கத்தக்கவாறு எண்ணிடத்தக்க
கணங்களே சாத்தியமான ``மிகக் கடினமான'' கணிக்கத்தக்கவாறு
எண்ணிடத்தக்க கணங்கள். \emph{எந்தவொரு} கணிக்கத்தக்கவாறு
எண்ணிடத்தக்க கணம் பற்றிய கேள்விகளுக்கும் விடையளிக்க அவை உதவுகின்றன.

\begin{thm}
$K$, $K_0$, $K_1$ அனைத்தும் முழுமையான கணிக்கத்தக்கவாறு
எண்ணிடத்தக்க கணங்கள்.
\end{thm}

\begin{proof}
$K_0$ முழுமையானது என்பதைக் காண, $B$ ஏதாவது ஒரு கணிக்கத்தக்கவாறு
எண்ணிடத்தக்க கணம் ஆகட்டும். அப்போது ஏதோ சுட்டெண் $e$ க்கு
\[
B = W_e = \Setabs{x}{\cfind{e}(x) \fdefined}.
\]
$f(x) = \tuple{e, x}$ என்ற சார்பாக $f$ ஐ எடுத்துக்கொள்க.
அப்போது ஒவ்வோர் இயல் எண் $x$ க்கும், $x \in B$ எனில், எனில்
மட்டுமே, $f(x) \in K_0$. வேறுவிதமாகக் கூறினால், $f$ ஆனது $B$
ஐ~$K_0$ க்குக் குறைக்கிறது.

$K_1$ முழுமையானது என்பதைக் காண, \olref[k1]{prop:k1} இன்
நிரூபணத்தில் $K_0$ ஐ அதற்குக் குறைத்தோம் என்பதைக் கவனிக்க.
எனவே \olref[ppr]{prop:trans-red} ஆல், எந்தக் கணிக்கத்தக்கவாறு
எண்ணிடத்தக்க கணத்தையும்~$K_1$ க்குக் குறைக்கலாம்.

இதே முறையில் $K$ ஐ $K_0$ க்குக் குறைக்கலாம்.
\end{proof}

\begin{prob}
$K$ இலிருந்து $K_0$ க்கான ஒரு குறைத்தலைத் தருக.
\end{prob}

\begin{digress}
இதுவரை நாம் கருதிய கணிக்கத்தக்கவாறு எண்ணிடத்தக்க கணங்களின் எல்லா
எடுத்துக்காட்டுகளும் கணிக்கத்தக்கவையாகவோ முழுமையானவையாகவோ
இருக்கின்றன. இது வியப்பாகத் தோன்ற வேண்டும்!{} கணிக்கத்தக்கதும்
அல்லாத, முழுமையானதும் அல்லாத கணிக்கத்தக்கவாறு எண்ணிடத்தக்க
கணங்களுக்கு எடுத்துக்காட்டுகள் உள்ளனவா? விடை ஆம்; ஆனால் 1950-களின்
நடுப்பகுதியில்தான் ஃபிரீட்பெர்க், முச்னிக் ஆகியோர் தனித்தனியாக
இதை நிறுவினர்.
\end{digress}

\end{document}

